\section{Boundaries of concentration and initialization}
\label{sec:boundaries}

The concentration theorem identifies the shape of a minimum-cost ready
state. Two separations delimit that statement. Nonlinear loss can make
distributed initial preparation strictly cheaper than every concentrated
state. Even under proportional loss, entering a sustainable ready regime
need not entail reaching any minimum-cost state in finite time.

\subsection{Smooth nonlinear loss can defeat concentration}
\label{sec:nonlinear-boundary}

Replace the proportional envelope by
$0\leq\rho_i(t)\leq g_i(p_i(t))$, keeping the other service, capacity,
and transfer assumptions unchanged. The next example uses identical
nonnegative, nondecreasing Lipschitz functions with $g_i(0)=0$.

\begin{proposition}[Nonlinear separation]
\label{prop:quadratic-failure}
For two identical modules with
\[
 M_1=M_2=1,\qquad a_1=a_2=\frac1{100},\qquad s=10,
 \qquad g_i(p)=p^2,\qquad H=\frac{23}{500},
\]
the ready state $(4/5,4/5)$ has maintenance cost $32/25$, while every
concentrated ready state has strictly larger cost.
\end{proposition}

\begin{proof}
Prepare module 1 at rate ten. Its loss is at most one, so its preparation
time is at most $1/45$. During that interval module 2 loses at most
$(16/25)/45$ and retains at least $884/1125$. After the first transfer,
capacity is $1001/100$ and the second module's net growth is at least
$901/100$. Its remaining deficit is at most $241/1125$. Hence the serial
policy's total time is at most
\[
 \frac1{45}+\frac{964}{40545}
 =\frac{373}{8109}<\frac{23}{500},
 \qquad
 \frac{23}{500}-\frac{373}{8109}=\frac7{4054500}>0.
\]
These are robust upper bounds, not exact quadratic traversal times.
Maintenance allocations $16/25$ to each coordinate preserve this target
under every allowed loss history.

For the converse, a concentrated state with no full module has total
initial preparation at most one. Transfers preserve the sum of completed
work and unfinished preparation, and nonnegative loss cannot increase it.
At least one additional unit of useful work is therefore necessary. Before
the last transfer, capacity is at most $1001/100$, whereas the deadline
permits at most
\[
 \frac{1001}{100}\frac{23}{500}
 =\frac{23023}{50000}<1
\]
units of allocated work. Such a state cannot be ready, under any schedule.

If one module is full and the other has preparation $r$, the same work
bound requires
\[
 r\geq 1-\frac{23023}{50000}=\frac{26977}{50000}.
\]
Its maintenance consequently satisfies
\[
 1+r^2-\frac{32}{25}
 \geq \frac{27758529}{2500000000}>0.
\]
Both modules full cost two. These cases exhaust concentrated states and
prove the claim.
\end{proof}

Thus smoothness, monotonicity, positive releases, and ample initial
capacity do not suffice for concentration. The proportional assumption is
substantive; equality of the proportional coefficients is not. The example
does not compute the nonlinear minimum itself.

\subsection{Critical readiness converges without necessarily becoming stationary}
\label{sec:startup-boundary}

We return to positive proportional coefficients. A normal warmup starts
from $p=0$, uses the same spare capacity $s$, and makes no transfers. A
removal deadline need only hold after the warmup. Write
$L(p)=\sum_i\gamma_i p_i$ and $Q(p)=\sum_i p_i$.

\begin{proposition}[Critical convergence]
\label{prop:critical-convergence}
Suppose $C(H)=s>0$. On the simultaneous maximal-loss history, every normal
trajectory that remains ready after a finite time $T_0$ converges to one
minimum-cost ready state $p^*$. Moreover,
\begin{equation}
 \int_{T_0}^{\infty}\bigl(L(p(t))-s+u(t)\bigr)\,dt
 \leq Q(T_0)-\lim_{t\to\infty}Q(t)<\infty.
 \label{eq:critical-dissipation}
\end{equation}
Finite-time arrival at $p^*$ is not asserted.
\end{proposition}

\begin{proof}
Reflection and the capacity constraint give
$Q'\leq s-u-L(p)$. Readiness implies $L(p)\geq C(H)=s$, so $Q$ is
nonincreasing and bounded below; integration proves
\eqref{eq:critical-dissipation}. Each coordinate is uniformly Lipschitz,
with derivative bounded in absolute value by
$\max\{s,\gamma_iM_i\}$. The nonnegative function $L(p(t))-s$ is therefore
uniformly Lipschitz and integrable, which forces it to tend to zero. A
positive lower bound along infinitely many times would otherwise produce
disjoint intervals of fixed positive area.

Compactness of the ready set implies that the distance to its minimizing
set tends to zero. That set is finite by Theorem~\ref{thm:concentration}:
the full set and partial identity leave the partial amount uniquely
determined by $L(p)=s$. A continuous trajectory eventually close to a
finite set cannot switch between disjoint neighborhoods of its members.
It therefore converges to one member.
\end{proof}

Convergence is not a reduction of startup to reaching a stationary
minimum. The following exact construction distinguishes those tasks.

\begin{theorem}[Startup without a reachable minimizing state]
\label{thm:startup-separation}
There is a three-module instance with positive sizes, coefficients,
releases, and spare capacity, and with finite cold exit, such that
$C(H)=s>0$ and its unique minimizing state cannot be robustly reached or
dominated after any finite normal warmup from cold. Nevertheless a finite normal
warmup enters a policy that is robustly $H$-ready at every later time.
\end{theorem}

\begin{proof}
Take $s=1$ and
\[
\begin{array}{c|ccc}
 \text{module}&M_i&\gamma_i&a_i\\ \hline
 A&1/2&1&2\\
 B&1&2&33\\
 C&3&1/2&1
\end{array}
\qquad
 p^*=(1/2,1/4,0),\qquad
 H=\tfrac12\log(5/2)+2\log(24/23).
\]
The target costs one. Transferring A immediately supplies rate three;
B then takes $\tfrac12\log(5/2)$. Its release supplies rate 36, so cold C
takes $2\log(24/23)$. Hence $p^*$ is ready.

We first certify that it is the unique minimum. The exact comparison
\begin{equation}
 H<\tfrac12\log3<\log2
 \label{eq:startup-deadline-comparison}
\end{equation}
follows from $5\cdot24^4=1658880<1679046=6\cdot23^4$.
Every minimizing state costs at most one and is concentrated by
Theorem~\ref{thm:concentration}. Neither B nor C can be full at that cost.
The remaining cases are exhaustive:
\begin{enumerate}
\item If A is subfull, it must complete first: subfull B and C cannot
complete at initial rate one, below their full loss rates two and $3/2$.
Cold A itself takes $\log2>H$. If A is partial, concentration makes B and C
cold; after A, their possible next stages take respectively
$\tfrac12\log3$ and $2\log2$, each exceeding $H$.
\item If A is full and B is the possible partial coordinate, cold C first
takes $2\log2>H$. B must precede C, giving total time
$\tfrac12\log(3-2p_B)+2\log(24/23)$. Readiness requires $p_B\geq1/4$,
while the cost bound $1/2+2p_B\leq1$ requires $p_B\leq1/4$.
\item If A is full and C is the possible partial coordinate, cost at most
one requires $p_C\leq1$. Cold B first already exceeds $H$. C first takes
at least $2\log(5/3)>\tfrac12\log3>H$, using $625>243$.
\end{enumerate}
Thus $C(H)=1$ and $p^*$ is its sole minimizer. The pure full/cold cases are
included by zero partial preparation. Seriality makes this an all-policy
certificate. Cold exit is finite in order A, B, C, with time
$\log2+\tfrac12\log3+2\log(24/23)$.

Next, no finite normal warmup can reach or dominate $p^*$ robustly.
On maximal loss, define $Z=p_A+p_B-3/4$. Every normal policy satisfies
\[
 Z'\leq1-p_A-2p_B=-2Z+(p_A-1/2)\leq-2Z.
\]
Starting cold gives $Z(t)\leq-(3/4)e^{-2t}<0$. Domination of $p^*$ would
instead require $Z\geq0$. Allocating to C cannot evade this support-budget
bound.

Finally, construct a different normal policy. Allocate $(0,0,1)$ for
$2\log2$, reaching the virtual maximal-loss state $(0,0,1)$. Thereafter
allocate $(1/2,1/2,0)$ forever. At elapsed time $t$ after the first phase,
put $x=e^{-t/2}$. The virtual state is
\begin{equation}
 p_A=\tfrac12(1-x^2),\qquad
 p_B=\tfrac14(1-x^4),\qquad p_C=x.
 \label{eq:startup-trajectory}
\end{equation}
For a request from this state, use order A, B, C. Its first two completion
times obey
\[
 e^{t_A}=1+x^2,\qquad
 e^{2t_{AB}}=Y(x):=\tfrac52+6x^2+\tfrac72x^4.
\]
The final completion time $T$ satisfies
\[
 e^{T/2}=\frac{72Y(x)^{1/4}-x}{69},\qquad
 e^{H/2}=\frac{72(5/2)^{1/4}}{69}.
\]
Concavity of the fourth root yields, for $0<x\leq1/128$,
\begin{align*}
 72\bigl(Y(x)^{1/4}-(5/2)^{1/4}\bigr)
 &\leq \frac{18}{(5/2)^{3/4}}\left(6x^2+\tfrac72x^4\right)\\
 &<108x^2+63x^4\\
 &\leq x\left(\frac{108}{128}+\frac{63}{128^3}\right)<x.
\end{align*}
Consequently $T<H$ for every finite request time after a total warmup of
$2\log2+14\log2=16\log2$. Every allocation uses the original normal
budget, and no module transfers before a request.

Under smaller allowed losses, scalar comparison keeps actual preparation
at least the virtual state. Readiness is upward closed, and the request
schedule can only improve. Thus the construction gives the claimed robust
guarantee, not only a maximal-loss trajectory calculation.
\end{proof}

Along \eqref{eq:startup-trajectory}, after readiness begins,
\[
 L(p)=1+\tfrac12x-\tfrac12x^2-\tfrac12x^4>1,\qquad
 Q(p)=\tfrac34+x-\tfrac12x^2-\tfrac14x^4\downarrow\tfrac34.
\]
All three coordinates remain partial at every finite time. The positive
preparation in C decays more slowly than the deficits in A and B, and its
request-time benefit offsets those deficits. Its total excess maintenance
is finite, consistent with Proposition~\ref{prop:critical-convergence}.
Concentration classifies minimizing states; it does not prohibit these
nonstationary ready states of strictly higher instantaneous maintenance.

The separation preserves the initialization distinction in the model. A
paid minimum is sufficient for the recurring frontier, while a finite
normal warmup may enter readiness through a different trajectory. Neither
construction supplies the removal guarantee before its stated entry time.
